\chapter{The Layered View: Rotations and Representations}
\label{ch:rotations}

\chapterorient{We have spent Part~III building a calculus---tensors and shapes,
the fold and its combinators, the state monad, operators as values. This closing
chapter steps back from the bricks to look at the building. The calculus is not
the simulator; it is one \emph{description} of the simulator, the topmost of
several. Below it sit the algebra, the loops, the pure functions, and finally the
cited C++ source. Each description is complete and correct in itself, and moving
from one to the next is a precise act we call a \emph{rotation}: a re-expression
that turns exactly one ``impedance''---mutation, fusion, iteration---while leaving
the computation untouched. By the end you will have the map of the whole project:
which layers exist, what each rotation buys, how to tell a genuine rotation from
an empty renaming, and---no less important---when a rotation honestly fails.}

\section{One simulator, several descriptions}

Throughout this book a single computation has worn many costumes. The
\texttt{axpy} update appeared first (\cref{ch:bigidea}) as a C loop overwriting an
array, then as a pure function returning a fresh tensor, then as an arrow in a
dependency graph. The transient solve appeared as a hand-written time-stepping
loop and again as a \code{fold\_solve} over a schedule. These were not different
computations. They were the \emph{same} computation, described at different
\emph{levels}.

That is the organizing idea of the entire project that this textbook documents.
We do not describe the simulator once. We describe it at a stack of levels, each a
faithful account of the machine, each written in a different vocabulary, and we
make the relationship between adjacent levels explicit and checkable. The
calculus you have just learned is the \emph{top} of that stack. This chapter
names all the levels, names the moves between them, and hands you the criteria
for telling a real move from a fake one.

\begin{keyidea}
The simulator is described at several \emph{layers}, $L_0$ through $L_4$, each
complete and correct \emph{in itself}. Adjacent layers are connected by a
\emph{rotation}: a re-expression that turns exactly one impedance---mutation,
fusion, iteration, or coordination---while computing the same numbers. The
calculus of Part~III is the top layer $L_4$; the physics of Part~II grounds the
bottom layers. The stack is the map; this chapter is the legend.
\end{keyidea}

\subsection{The five layers}

We number the layers from the ground up. The bottom is the literal source; the
top is the calculus. \Cref{fig:stack} is the hero picture of the whole chapter,
and indeed of the project---print it on the inside of your eyelids.

\begin{description}[leftmargin=2.2em,style=nextline]
  \item[$L_0$ --- the cited source.] The ground truth: ranges of Palace's
  C++ (and the MFEM library beneath it), quoted by file and line. Every claim
  anywhere in the stack ultimately rests on an $L_0$ citation. $L_0$ is not
  ``ours''; it is the machine as written. Nothing is interpreted here---it is read.

  \item[$L_1$ --- pure functions.] The same operations re-expressed as
  \emph{pure functions over immutable values}. The destination array that $L_0$
  overwrote in place becomes a returned value; the loop that scribbled over memory
  becomes a function from inputs to a fresh output. This is the
  \emph{mutation rotation} of \cref{ch:bigidea}, applied operator by operator.

  \item[$L_2$ --- algebraic decompositions.] Each $L_1$ function unfolded back
  into a \emph{composition of base algebraic primitives}, with the performance
  tricks erased. A fused kernel that computed $a x + b y + c z$ in one
  cache-friendly sweep is decomposed into the scaled sums it stands for; a
  packed, blocked matrix apply becomes the plain linear map it implements. This is
  the \emph{fusion rotation}.

  \item[$L_3$ --- global tensor-field operations.] Where the algebra permits,
  the explicit per-degree-of-freedom loops of $L_2$ are lifted to
  \emph{whole-field operations}: ``add these two fields,'' ``apply this operator
  to this field,'' with no surviving element loop in the spec form. This is the
  \emph{iteration rotation}---and, crucially, the one that does not always go
  through (\cref{sec:obstruction}).

  \item[$L_4$ --- the calculus.] The small, formally defined
  graph-evaluation calculus of Part~III: high-order combinators
  (\cref{ch:fold,ch:combinators}), the state monad (\cref{ch:monad}), the strata
  (\cref{ch:strata}), operators as values (\cref{ch:operators}). $L_4$ is
  \emph{vocabulary, not architecture}---it says what the algorithm \emph{is},
  abstractly, in a form an accelerator backend can consume. The
  \emph{coordination rotation} from $L_3$ replaces hand-written control flow with
  named combinators.
\end{description}

\begin{figure}[t]
\centering
\begin{tikzpicture}[node distance=0mm]
  % five stacked layer bands, L0 at bottom -> L4 at top
  \def\lw{86mm}   % layer width
  \def\lh{11mm}   % layer height
  \def\gap{8.5mm} % vertical gap between bands (room for rotation arrow)
  \foreach \i/\name/\desc/\col/\bg in {%
      0/$L_0$/cited C++ {\&} MFEM source/simInk/simBgGray,
      1/$L_1$/pure functions over immutable values/simGreen/simBgGreen,
      2/$L_2$/algebraic decompositions {\textbullet} tricks erased/simTeal/simBgGray,
      3/$L_3$/global tensor-field operations/simBlue/simBgBlue,
      4/$L_4$/calculus {\textbullet} combinators {\textbullet} monad {\textbullet} strata/simViolet/simBgViolet}{
    \node[draw=\col, fill=\bg, thick, rounded corners=2pt,
          minimum width=\lw, minimum height=\lh, align=center, font=\small]
      (L\i) at (0, \i*\dimexpr\lh+\gap\relax) {%
        \textbf{\name}\quad\textcolor{simGray}{\footnotesize\desc}};
  }
  % rotation arrows between adjacent bands
  \foreach \lo/\hi in {0/1,1/2,2/3,3/4}{
    \draw[depedge, line width=1pt] (L\lo.north) -- (L\hi.south);
  }
  % labels placed to the right of each arrow gap
  \node[font=\scriptsize\itshape, simBlue, anchor=west] at (0.50*\lw, 0.5*\lh+0.5*\gap) {mutation rotation};
  \node[font=\scriptsize, simGray, anchor=west] at (0.50*\lw, 0.5*\lh+0.5*\gap-3mm) {mutation $\to$ purity};
  \node[font=\scriptsize\itshape, simBlue, anchor=west] at (0.50*\lw, 1.5*\lh+1.5*\gap) {fusion rotation};
  \node[font=\scriptsize, simGray, anchor=west] at (0.50*\lw, 1.5*\lh+1.5*\gap-3mm) {fused kernel $\to$ algebra};
  \node[font=\scriptsize\itshape, simBlue, anchor=west] at (0.50*\lw, 2.5*\lh+2.5*\gap) {iteration rotation};
  \node[font=\scriptsize, simGray, anchor=west] at (0.50*\lw, 2.5*\lh+2.5*\gap-3mm) {loop $\to$ whole-field op};
  \node[font=\scriptsize\itshape, simBlue, anchor=west] at (0.50*\lw, 3.5*\lh+3.5*\gap) {coordination rotation};
  \node[font=\scriptsize, simGray, anchor=west] at (0.50*\lw, 3.5*\lh+3.5*\gap-3mm) {control flow $\to$ combinators};
  % side bracket: "complete & correct in itself"
  \draw[decorate, decoration={brace, amplitude=5pt}, simGray]
    (-0.52*\lw-2mm, -0.5*\lh) -- (-0.52*\lw-2mm, 4*\dimexpr\lh+\gap\relax+0.5*\lh);
  \node[rotate=90, font=\scriptsize\itshape, simGray, anchor=south]
    at (-0.52*\lw-7mm, 2*\dimexpr\lh+\gap\relax) {each layer complete \& correct in itself};
\end{tikzpicture}
\caption[The five-layer stack and its rotations]{The impedance-matching stack.
Read it from the ground up: $L_0$ is the cited source, and each rotation above it
turns exactly one impedance while computing the same numbers---mutation into
purity ($L_0\!\to\!L_1$), a fused kernel into its base algebra
($L_1\!\to\!L_2$), an explicit loop into a whole-field operation
($L_2\!\to\!L_3$), and hand-written control flow into named combinators
($L_3\!\to\!L_4$). Every layer is a complete, correct account of the simulator in
its own vocabulary; the arrows are translations between vocabularies, never mere
renamings. The book has been teaching the $L_4$ view, with Part~II supplying the
$L_0$--$L_1$ grounding.}
\label{fig:stack}
\end{figure}

\subsection{Why a stack and not a single description}

It would be simpler to describe the simulator once and have done. The stack earns
its complexity because \emph{each layer answers a different question well, and
none answers all of them}. $L_0$ answers ``what does Palace actually do''---it is
the only layer that can be wrong about the source, because it \emph{is} the
source. $L_4$ answers ``what algorithm is this, abstractly, and how would it lower
onto a GPU''---but it is too far from the source to verify against it directly.
The intervening layers are the stepping-stones that let us connect the two with
small, individually checkable moves. A direct $L_0 \to L_4$ leap would be a leap
of faith; the stack replaces it with a sequence of short, auditable steps.

This is why we call the stack \emph{impedance-matching}. In an electrical
transmission line, you do not connect a high-impedance source straight to a
low-impedance load---you insert matching sections that each step the impedance part
of the way, so power transfers without reflection. Here the ``impedance mismatch''
is between raw mutating C++ and an abstract whole-tensor calculus. Each rotation
matches one stage of it.

\section{What a rotation is}
\label{sec:whatis}

We have used the word loosely; now we pin it down. A rotation is the unit of
progress between layers. The metaphor is geometric: a rotation changes your
\emph{point of view} on an object without changing the object. Rotate a cube and
the face that was hidden comes forward, but it is the same cube, with the same
volume and the same corners.

\begin{definition}[Rotation]
\label{def:rotation}
A \emph{rotation} is a re-expression of a piece of work from layer $L_n$ into
layer $L_{n+1}$ that \emph{changes exactly one impedance} while computing the same
result. The impedance turned is one of: mutation (made into purity), a fused
kernel (unfolded into base algebra), an explicit loop (lifted to a whole-field
operation), or hand-written coordination (replaced by a named combinator).
\end{definition}

The single most important thing to understand about a rotation is what it is
\emph{not}.

\begin{keyidea}
A rotation is a \emph{translation across vocabularies}, not a one-to-one
\emph{renaming}. It re-expresses the lower form in a genuinely different idiom,
in which something becomes simpler, hidden, or substitutable. If the new form is
just the old form with the identifiers swapped---same state threaded, same
structure, same grain---no impedance turned, and \emph{no rotation occurred}. A
rotation turns one impedance; a renaming turns none.
\end{keyidea}

To make this concrete, contrast two purported $L_1 \to L_2$ moves on the
conjugate-gradient inner step. At $L_1$, one step reads
\[
  \vr \leftarrow \vr - \alpha\,\mat{A}\vp, \qquad
  \beta \leftarrow \tfrac{(\vr\cdot\vr)_{\text{new}}}{(\vr\cdot\vr)_{\text{old}}},\qquad
  \vp \leftarrow \vr + \beta\,\vp .
\]
A \emph{renaming} rewrites this as \lstinline[style=l4style]|axpby(alpha, Ap, -1, r)|;
\lstinline[style=l4style]|beta = dot(r,r)/dot(r_old,r_old)|;
\lstinline[style=l4style]|axpby(1, r, beta, p)|. Every line is the same operation
with a BLAS-1 name pasted on. The reader could not drop in a \emph{different
algorithm} where \code{axpby} sits---there is only the one operation it stands
for. Nothing was hidden, nothing compressed, nothing made substitutable. This is
the failure mode \cref{sec:tests} teaches you to detect.

A genuine rotation, by contrast, introduces a primitive
\lstinline[style=l4style]|cg_step| whose signature \emph{hides} the explicit
threading of $\vr_k \to \vr_{k+1}$ and into whose slot a \emph{different}
inner-step algorithm---\code{cgne\_step}, \code{minres\_step}---could be
substituted without the caller noticing. That is a translation into a new
vocabulary, not a renaming in the old one.

\section{The three tests for ``a rotation has occurred''}
\label{sec:tests}

How do you \emph{tell}? This is the analytical core of the chapter, and the one
piece of it you should carry out the door. A proposed rotation from $L_n$ to
$L_{n+1}$ is genuine if and only if \textbf{at least one} of three tests passes.
If none passes, you have a renaming---a smell, not a rotation. \Cref{fig:tests}
collects the tests as a decision diagram.

\begin{figure}[t]
\centering
\begin{tikzpicture}[node distance=7mm]
  \node[data, text width=30mm] (in) {a proposed $L_n\!\to\!L_{n+1}$ move};
  \node[stage, below=of in, text width=46mm, minimum height=10mm] (t1)
    {\textbf{T1.} Does it \emph{hide state}?\\[1pt]\footnotesize something mutable now internal/absorbed};
  \node[stage, below=8mm of t1, text width=46mm, minimum height=10mm] (t2)
    {\textbf{T2.} Does a \emph{different algorithm}\\fit the same slot?\\[1pt]\footnotesize the algorithmic-substitution test};
  \node[stage, below=8mm of t2, text width=46mm, minimum height=10mm] (t3)
    {\textbf{T3.} Is the \emph{threaded state}\\compressed or abstracted?};
  % outcome nodes on the right
  \node[product, right=20mm of t1, text width=24mm] (yes) {genuine\\rotation};
  \node[draw=simRust, fill=simBgRust, thick, rounded corners=2pt, align=center,
        font=\small\bfseries, text=simInk, right=20mm of t3, text width=24mm,
        minimum height=9mm] (no) {renaming\\(smell)};
  % flow
  \draw[flow] (in) -- (t1);
  \draw[flow] (t1) -- node[left,font=\scriptsize]{no} (t2);
  \draw[flow] (t2) -- node[left,font=\scriptsize]{no} (t3);
  \draw[flow] (t1.east) -- node[above,font=\scriptsize]{yes} (yes.west);
  \draw[flow] (t2.east) -| node[pos=0.25,above,font=\scriptsize]{yes} (yes.south);
  \draw[flow] (t3.east) -- node[above,font=\scriptsize]{yes} ($(t3.east)+(8mm,0)$) |- (yes.south west);
  \draw[flow] (t3.south) |- node[pos=0.75,below,font=\scriptsize]{no to all three} (no.west);
\end{tikzpicture}
\caption[The three rotation tests as a decision diagram]{The three tests for a
genuine rotation. A move passes if \emph{any one} of them is satisfied: it hides
mutable state (\textbf{T1}), it admits a different algorithm in the same slot
(\textbf{T2}, the algorithmic-substitution test), or it compresses the state
threaded through the computation (\textbf{T3}). Only when \emph{all three} fail is
the move a mere renaming---the smell this chapter teaches you to catch.}
\label{fig:tests}
\end{figure}

\subsection{T1 --- state hiding}

\emph{Does the higher form hide a piece of state that the lower form exposed?}

This is the most common way a rotation succeeds. At $L_n$ some buffer, index,
accumulator, or configuration is visible---threaded through the code, named in
signatures, mutated by hand. At $L_{n+1}$ it is \emph{absorbed}: tucked inside a
constructed operator, captured at construction time, or hidden behind a
primitive's contract.

The canonical example is a constructed operator (\cref{ch:operators}). A
preconditioner is configured once---factorizations computed, levels assembled,
tables filled---and then applied many times. At a low layer all that
configuration is visible state, threaded into every apply. The rotation
\emph{constructs} an operator that captures the configuration internally and
exposes only \lstinline[style=l4style]|apply :: LinOp -> Tensor[N] -> Tensor[N]|.
The construction-time state is hidden; the caller sees a clean field-to-field map.

\subsection{T2 --- coarser substitution (the algorithmic-substitution test)}

\emph{Could you swap the implementation in this slot for a genuinely different
algorithm without the caller noticing?}

This is the subtlest and most valuable test, so we give it its own name: the
\emph{algorithmic-substitution test}. The trap it guards against is mistaking a
\emph{name} for an \emph{interface}. Pasting the name \code{dot} onto an inner
product does not make a coarser substitution, because there is only one algorithm
called \code{dot}. The substitution must be \emph{algorithmic}, not nominal.

\begin{notationbox}
The algorithmic-substitution test, stated as a question to ask before claiming a
rotation: \emph{could a reader replace the $L_{n+1}$ primitive with a different
algorithm---not a different implementation of the same algorithm---and still
satisfy the $L_n$ contract?} If yes, a coarser substitution interface emerged and
the rotation is genuine. If ``no, it is just a named version of the one
operation,'' it is a renaming.
\end{notationbox}

A passing example: at $L_n$, choosing modified Gram--Schmidt over classical
Gram--Schmidt means re-threading per-vector dot accumulators and rewriting the
inner loop. At a genuine $L_{n+1}$ the choice collapses to a single substitution,
\lstinline[style=l4style]@orthog := mgs_step | cgs_step | cgs2_step@: three
\emph{distinct algorithms} occupying one primitive's slot, with whole-algorithm
correctness arguments factoring through the primitive's contract rather than its
body. \Cref{fig:substitution} draws the shape of a passing T2: two unlike
implementations, one interface.

\begin{figure}[t]
\centering
\begin{tikzpicture}[node distance=8mm]
  % the interface
  \node[opbox, minimum width=34mm, minimum height=9mm] (iface)
    {\lstinline[style=l4style]|apply :: LinOp -> Tensor[N] -> Tensor[N]|};
  \node[font=\scriptsize\itshape, simGray, above=1mm of iface] {one interface (the slot)};
  % caller above
  \node[data, above=9mm of iface, text width=26mm] (caller) {the caller\\(can't tell which)};
  \draw[flow] (caller) -- node[right,font=\scriptsize]{\code{apply}} (iface);
  % two unlike implementations below
  \node[extern, below left=11mm and 2mm of iface, text width=30mm, minimum height=11mm] (impA)
    {\textbf{Jacobi smoother}\\[1pt]\footnotesize diagonal scaling, one sweep};
  \node[extern, below right=11mm and 2mm of iface, text width=30mm, minimum height=11mm] (impB)
    {\textbf{multigrid V-cycle}\\[1pt]\footnotesize restrict, recurse, prolong};
  \draw[flow, densely dashed, simViolet] (iface.south) -- (impA.north);
  \draw[flow, densely dashed, simViolet] (iface.south) -- (impB.north);
  \node[font=\scriptsize\itshape, simGray, below=1mm of impA] {algorithm 1};
  \node[font=\scriptsize\itshape, simGray, below=1mm of impB] {algorithm 2};
  \node[align=center, font=\scriptsize\itshape, simGray, below=15mm of iface]
    {two \emph{different} algorithms, one slot $\Rightarrow$ T2 passes};
\end{tikzpicture}
\caption[The algorithmic-substitution test]{The shape of a passing
algorithmic-substitution test (T2). A single interface---here a preconditioner's
\code{apply}---sits between the caller and the implementation. Behind it can stand
genuinely \emph{different} algorithms: a one-line Jacobi diagonal scaling, or a
full multigrid V-cycle that restricts, recurses, and prolongs. The caller invokes
\code{apply} and cannot tell which is in the slot. Because a different algorithm
fits the same slot, the move is a genuine rotation. Contrast a renaming, where the
``slot'' admits only the one operation its name stands for.}
\label{fig:substitution}
\end{figure}

\subsection{T3 --- threaded-state compression}

\emph{Is the bundle of state threaded through each step strictly smaller, or at
least strictly more abstract, at the higher layer?}

A solver's inner loop may thread a long tuple of buffers: a Krylov basis array,
a packed Hessenberg matrix, Givens rotation accumulators, a column index, scratch
vectors, a restart counter. If the higher layer threads a short, abstract bundle
instead---\lstinline[style=l4style]|(iterate, residual, basis_handle, converged)|,
with the implementation buffers encapsulated behind \code{basis\_handle}---then
the threaded state compressed, and the rotation is genuine even if no single
piece was outright hidden and no algorithm became substitutable. The reader
manipulates fewer, more meaningful names.

\begin{pitfall}
The renaming smell is the trap to watch for: a proposed higher form that names
the lower form's operations as primitives but turns \emph{no} impedance. Symptom:
the higher form threads the \emph{same} buffers with the \emph{same} shapes,
admits only the \emph{same} algorithm in each slot, and exposes the \emph{same}
state. A genuine rotation may pass \emph{some} pieces through unchanged---an
\code{axpy} that is already at the right grain need not be ``rotated'' for its own
sake---but the anti-pattern is \emph{everything} carrying through with new names
and nothing moving forward. When you find it, do not record a hollow layer:
either merge (treat the lower form as the highest this slice has reached) or
redesign the move so one of the three tests genuinely fires.
\end{pitfall}

\section{Justification: how we know a rotation is sound}
\label{sec:justification}

Detecting that a rotation \emph{occurred} (\cref{sec:tests}) is separate from
showing it is \emph{correct}---that the higher form really computes what the lower
form does. Every rotation in the stack carries a \emph{justification}, and there
are exactly five kinds. \Cref{fig:justification} is the palette. The first four
are positive (the rotation goes through, here is why); the fifth is the negative
result that the next section is about.

\begin{description}[leftmargin=2.2em,style=nextline]
  \item[\texttt{algebraic}.] An algebraic identity makes the rotation manifest:
  associativity, distributivity, linearity. Unfolding a fused $a x + b y$ into two
  scaled sums is justified by the distributive law---no test required, the
  identity \emph{is} the proof.

  \item[\texttt{structural}.] A structural property of the primitives carries the
  rotation: the locality of an element-wise operation, the commutativity of an
  operator apply with a scalar scaling. The reasoning is about \emph{shape}, not
  arithmetic value.

  \item[\texttt{reduction\_chain}.] A sequence of small algebraic or structural
  steps, each individually named and sound, composed into the whole rotation. When
  no single identity does the job, a documented chain of them does.

  \item[\texttt{empirical\_match}.] An \emph{executed test} confirms the higher
  form matches the lower form numerically: run both on the same input, check the
  numbers agree to tolerance. \textbf{This is the preferred justification whenever
  it is available}---a green test beats a paragraph of algebra, because it cannot
  quietly be wrong about an edge case the prose forgot.

  \item[\texttt{obstruction}.] The rotation does \emph{not} go through, and the
  justification records \emph{why} as a first-class result. This is
  \cref{sec:obstruction}.
\end{description}

\begin{figure}[t]
\centering
\begin{tikzpicture}[node distance=4mm]
  \def\pw{30mm}
  \node[draw=simGreen, fill=simBgGreen, thick, rounded corners=2pt, align=center,
        font=\footnotesize, text=simInk, minimum height=13mm, text width=\pw] (alg)
    {\textbf{\code{algebraic}}\\[1pt]an identity makes it manifest};
  \node[draw=simTeal, fill=simBgGray, thick, rounded corners=2pt, align=center,
        font=\footnotesize, text=simInk, minimum height=13mm, text width=\pw,
        right=of alg] (str)
    {\textbf{\code{structural}}\\[1pt]a property of the shape carries it};
  \node[draw=simBlue, fill=simBgBlue, thick, rounded corners=2pt, align=center,
        font=\footnotesize, text=simInk, minimum height=13mm, text width=\pw,
        right=of str] (red)
    {\textbf{\code{reduction\_chain}}\\[1pt]named small steps, composed};
  \node[draw=simGold, fill=simBgGold, very thick, rounded corners=2pt, align=center,
        font=\footnotesize, text=simInk, minimum height=13mm, text width=\pw,
        below=6mm of str] (emp)
    {\textbf{\code{empirical\_match}}\\[1pt]\emph{run both; numbers agree}\\(preferred)};
  \node[draw=simRust, fill=simBgRust, very thick, rounded corners=2pt, align=center,
        font=\footnotesize, text=simInk, minimum height=13mm, text width=\pw,
        right=of emp] (obs)
    {\textbf{\code{obstruction}}\\[1pt]it does \emph{not} rotate; record why};
  \node[font=\scriptsize\itshape, simGray, left=of emp, text width=24mm]
    {a first-class\\negative result $\rightarrow$};
\end{tikzpicture}
\caption[The justification-kinds palette]{The five kinds of justification a
rotation can carry. The four upper/positive kinds certify that the rotation goes
through: an algebraic identity, a structural property, a chain of small named
steps, or---preferred whenever feasible---an empirical match, where the two layers
are run on the same input and their numbers are checked to agree. The fifth,
\code{obstruction} (rust), is the first-class \emph{negative} result: the rotation
genuinely does not go through, and recording \emph{why} is itself valuable
output.}
\label{fig:justification}
\end{figure}

\section{Obstructions: when a rotation honestly fails}
\label{sec:obstruction}

Not everything rotates. Some work, faithfully described at $L_2$, genuinely cannot
be lifted to a clean $L_3$ whole-field operation, because its very meaning depends
on doing things in order. When that happens we do \emph{not} fake a rotation and
we do \emph{not} treat the failure as an error in our method. We record an
\emph{obstruction}: a precise statement of what blocks the lift and why. An
obstruction is a finding, not a failure.

\begin{keyidea}
An \emph{obstruction} is a first-class negative result: a rotation that genuinely
does not go through, with the reason recorded as part of the spec. The obstructions
are not gaps in the work---they are some of its most valuable output, because they
say exactly where the abstract whole-tensor picture meets a wall, and why. ``It
does not rotate, and here is precisely why'' is knowledge.
\end{keyidea}

The archetype is the \emph{sequential reduction}---a loop whose $i$-th step reads
the $(i-1)$-th step's \emph{output}. \Cref{ch:tensors} introduced exactly this as
the case where the tensor-field lift fails: a Gauss--Seidel sweep updates
component $i$ using the \emph{already-updated} components $0,\dots,i-1$, so there
is no way to apply it to ``all entries at once.'' A whole-field operation has no
notion of ``the entry before this one''; the loop-carried dependency is the
obstruction. \Cref{fig:obstruction} contrasts a liftable elementwise update with
an unliftable sequential one.

\begin{figure}[t]
\centering
\begin{tikzpicture}[node distance=6mm]
  % ---- LEFT: liftable, independent ----
  \begin{scope}
    \node[font=\bfseries\small, simGreen] at (0,1.7) {Liftable: independent updates};
    \foreach \i in {0,1,2,3}{
      \node[opbox, minimum width=8mm, minimum height=6mm] (a\i) at (\i*0.95,0.7) {$y_\i$};
      \node[data, minimum width=8mm, minimum height=5mm] (x\i) at (\i*0.95,-0.4) {$x_\i$};
      \draw[flow] (x\i) -- (a\i);
    }
    \node[align=center, font=\scriptsize\itshape, simGray, text width=42mm] at (1.4,-1.25)
      {each $y_i = a x_i + y_i$ reads only its own input $\Rightarrow$ apply to all at once};
    \node[draw=simGreen, thick, rounded corners=2pt, fit=(a0)(a3), inner sep=3pt,
          fill=simGreen!8] {};
    \node[font=\scriptsize\bfseries, simGreen] at (1.4,1.15) {one whole-field op};
  \end{scope}
  % ---- RIGHT: unliftable, sequential ----
  \begin{scope}[shift={(6.3,0)}]
    \node[font=\bfseries\small, simRust] at (0,1.7) {Unliftable: sequential reduction};
    \foreach \i in {0,1,2,3}{
      \node[opbox, draw=simRust, minimum width=8mm, minimum height=6mm] (b\i) at (\i*0.95,0.7) {$y_\i$};
      \node[data, minimum width=8mm, minimum height=5mm] (u\i) at (\i*0.95,-0.4) {$x_\i$};
      \draw[flow] (u\i) -- (b\i);
    }
    % chain of carried dependencies
    \foreach \i/\j in {0/1,1/2,2/3}{
      \draw[backflow] (b\i.east) -- (b\j.west);}
    \node[align=center, font=\scriptsize\itshape, simGray, text width=42mm] at (1.4,-1.25)
      {$y_i$ reads the \emph{already-updated} $y_{i-1}$ $\Rightarrow$ no whole-field form};
    \node[font=\scriptsize\bfseries, simRust] at (1.4,1.15) {loop-carried dependency};
  \end{scope}
\end{tikzpicture}
\caption[An obstruction: a sequential reduction will not lift]{Why some loops
rotate to $L_3$ and some do not. \emph{Left:} an \code{axpy}-style update writes
$y_i = a x_i + y_i$ reading only its own input; the entries are independent, so the
loop lifts to a single whole-field operation. \emph{Right:} a sequential
reduction (a Gauss--Seidel sweep) writes $y_i$ using the \emph{already-updated}
$y_{i-1}$; the dashed rust arrows are the loop-carried dependency. A whole-field
operation has no ``previous entry,'' so the lift fails---and we record the
obstruction (with the offending dependency cited) rather than force a false
rotation. This is the tensor-field-lift failure of \cref{ch:tensors}.}
\label{fig:obstruction}
\end{figure}

A subtle and important variant: sometimes the loop we cannot rotate is one
\emph{we never wrote}. The eigenmode analysis (\cref{ch:targets}) hands its operators
to an opaque library routine---SLEPc's eigensolver---whose iteration lives entirely
in someone else's code. We cannot render that loop as a fold because Palace authors
no loop to render; we can only mark the obstruction at the library boundary. This
is still a genuine, recorded obstruction; it just sits at a wrapper, not at a
visible recurrence. We meet exactly this kind again in \cref{ch:eigenvalues}.

\begin{remark}
An obstruction is never an excuse to stop. It is a precise boundary marker: the
work \emph{below} it (the per-step body) usually rotates cleanly, and only the
\emph{loop} resists. The honest record says ``the body lifts; the iteration does
not, because of this cited dependency''---which is far more useful than either a
forced rotation or a blank silence.
\end{remark}

\section{Two worked examples}

We now run the machinery on two concrete cases: a clean state-hiding rotation, and
a clean algorithmic-substitution rotation set beside a renaming that fails all
three tests.

\begin{workedexample}{The three tests on the \texttt{axpy} mutation rotation}{axpy}
Take the $L_0 \to L_1$ \emph{mutation rotation} on \code{axpy}, the running
example of \cref{ch:bigidea}. At $L_0$, the cited C++ overwrites the destination
in place:
\begin{lstlisting}[style=l4style, language=]
for (int i = 0; i < n; i++)
    y[i] = a * x[i] + y[i];     // y is mutated; the name y survives, values do not
\end{lstlisting}
At $L_1$, the operation is a pure function returning a fresh tensor:
\begin{lstlisting}[style=l4style]
axpy :: Scalar -> Tensor[N] -> Tensor[N] -> Tensor[N]
axpy a x y = a * x + y         -- returns a FRESH tensor; x and y unchanged
\end{lstlisting}

Apply the three tests of \cref{sec:tests}.

\textbf{T1 (state hiding): passes.} At $L_0$ the destination array \code{y} is a
piece of \emph{mutable state}---a block of memory the loop scribbles over, whose
contents depend on how many iterations have run. At $L_1$ that mutable destination
is gone: \code{axpy} \emph{returns} its result, and its inputs \code{x} and
\code{y} are untouched afterward. The mutable destination-argument has been hidden
(in fact, eliminated): the impedance \emph{mutation} was turned into \emph{purity}.
That is precisely the $L_0\!\to\!L_1$ impedance of \cref{def:rotation}, and it is
the one test the mutation rotation is built to pass.

\textbf{T2 (coarser substitution): does not pass---and need not.} Could a reader
drop a \emph{different algorithm} into the \code{axpy} slot and still satisfy the
$L_0$ contract? No: \code{axpy} computes $a x + y$ and there is exactly one such
operation. The slot admits no algorithmic choice. T2 fails---but a rotation needs
only \emph{one} test to pass, and T1 already did. \code{axpy} is a leaf primitive,
already at the right grain; it is not where substitution interfaces live.

\textbf{T3 (threaded-state compression): not applicable.} \code{axpy} is a single
operation, not an iteration threading a state bundle; there is no per-step state to
compress. T3 has nothing to bite on here.

\textbf{Verdict.} \emph{One} test passes (T1, decisively), so the move is a
genuine rotation, not a renaming. Its justification (\cref{sec:justification}) is
\code{algebraic}: the equality of the in-place result and the returned value is the
identity $a x + y = a x + y$ read under two memory models---and it is also trivially
confirmable by \code{empirical\_match}, running both forms on the same \code{x},
\code{y}, \code{a} and checking the output array agrees entry for entry.
\end{workedexample}

\begin{workedexample}{Solver-as-operator: a passing T2 beside a failing renaming}{slot}
Now the algorithmic-substitution test in action, on a \emph{preconditioner behind
an apply interface}---the ``solver-as-operator'' pattern of \cref{ch:operators}.
A preconditioner $\mat{M}^{-1}$ is used inside a Krylov solve only through one
move: \emph{apply it to a vector}. The rotation packages it as a constructed
operator exposing a single method:
\begin{lstlisting}[style=l4style]
apply :: LinOp -> Tensor[N] -> Tensor[N]      -- the only thing the solver calls
\end{lstlisting}

\textbf{The passing rotation (T2).} Ask the algorithmic-substitution question:
could a \emph{different algorithm} occupy this \code{apply} slot while the Krylov
solver---the caller---stays byte-for-byte the same? Yes, and dramatically so. The
slot can hold a one-line Jacobi smoother (scale by the inverse diagonal); or a
Gauss--Seidel sweep; or a full geometric-multigrid V-cycle that restricts the
residual to a coarse mesh, recurses, and prolongs the correction back
(\cref{ch:multigrid}). These are \emph{wildly} different algorithms---different
work, different cost, different code---yet each satisfies the same contract
``approximately invert the operator,'' and the caller, which only ever writes
\lstinline[style=l4style]|z = apply(pc, r)|, cannot tell which is installed
(\cref{fig:substitution}). A different algorithm fits the same slot: T2 passes, and
this is a genuine rotation. (T1 also passes here---the preconditioner's setup state
is hidden inside the constructed operator---which is common: real rotations often
trip more than one test.)

\textbf{The failing renaming, for contrast.} Suppose instead we had ``rotated''
the diagonal-scaling preconditioner by simply renaming its inner expression
$z_i \leftarrow r_i / d_i$ to a primitive \code{diag\_scale} and calling it a new
layer:
\begin{lstlisting}[style=l4style]
diag_scale :: Tensor[N] -> Tensor[N] -> Tensor[N]   -- z[i] = r[i] / d[i]
\end{lstlisting}
Run the tests. \textbf{T1}: nothing is hidden---the diagonal \code{d} is still
threaded in as an explicit argument, just as before. \textbf{T2}: no different
algorithm fits the \code{diag\_scale} slot; it \emph{is} the one elementwise
division, named. \textbf{T3}: there is no threaded state bundle to compress; it is
a single stateless operation. All three fail. This is a renaming---the smell of
\cref{sec:tests}---and the right response is not to enshrine it as a layer but to
recognize \code{diag\_scale} as a leaf that simply \emph{carries through}, while
the real rotation happens one level up, at the \code{apply} interface that makes
the \emph{whole preconditioner} substitutable.

\textbf{The lesson.} The same component---a diagonal preconditioner---hosts both a
genuine rotation and an empty one. The rotation is genuine at the grain where
\emph{substitution becomes possible} (the \code{apply} interface, where Jacobi and
multigrid are interchangeable) and empty at the grain where it does not (renaming
one elementwise division). The three tests are what tell the two apart.
\end{workedexample}

\section{An advanced sidebar: erasure scope}
\label{sec:erasure}

\emph{This section is a depth note for the curious; a first reading may skip to
\cref{sec:payoff}.}

When an iteration rotation \emph{does} engage an obstruction
(\cref{sec:obstruction}), there is a finer question worth asking: \emph{how much}
of the operator's iteration view does the layer-hop erase? The answer ranges over
a small ladder---the \emph{erasure scope}---and naming the rung sharpens what each
obstruction is really about. \Cref{fig:erasure} draws the ladder.

\begin{figure}[t]
\centering
\begin{tikzpicture}[node distance=5mm]
  \node[band, fill=simBgGreen, draw=simGreen, thick, text width=68mm] (r1)
    {\textbf{1. single loop}\;---\;the whole operator \emph{is} one sequential loop
     (a linear solve's outer iteration). One recurrence erased.};
  \node[band, fill=simBgGray, draw=simTeal, thick, text width=68mm, below=of r1] (r2)
    {\textbf{2. nested double loop}\;---\;an inner recurrence inside an outer sweep,
     both sequential (a polynomial smoother). Two nested recurrences erased.};
  \node[band, fill=simBgRust, draw=simRust, thick, text width=68mm, below=of r2] (r3)
    {\textbf{3. opaque library call}\;---\;the loop lives entirely in
     someone else's code (an eigensolver). No loop to render: only a boundary
     marker erased.};
  % ascending "more iteration erased" arrow on the left
  \draw[-{Stealth[length=2.5mm]}, simGray, thick] (-3.95,-2.2) -- (-3.95,1.0);
  \node[font=\scriptsize\itshape, simGray, rotate=90, anchor=south]
    at (-4.3,-0.6) {more of the iteration view erased};
\end{tikzpicture}
\caption[The erasure-scope ladder]{How much iteration view a layer-hop erases,
from least to most opaque. \emph{Rung 1:} the operator is a single sequential
loop---one linear-solve outer iteration---and the hop erases that one recurrence.
\emph{Rung 2:} a nested double loop, an inner recurrence inside an outer sweep (a
degree-$k$ polynomial smoother), erasing two coupled recurrences. \emph{Rung 3:}
an opaque library call---an eigensolver whose loop is in third-party code---where
there is no loop to render at all, only a boundary marker to erase. The rung tells
you what kind of obstruction you are looking at and how it must be recorded.}
\label{fig:erasure}
\end{figure}

The three rungs differ in a structurally meaningful way. On rungs~1 and~2 Palace
\emph{authors} the recurrence, so the lower layer can \emph{render} it---an
explicit tail-recursive loop carrying a visible loop-carried dependency---and the
hop then \emph{erases} that rendered iteration view, the obstruction surviving only
as ``these algebraic laws do \emph{not} hold'' (you cannot reorder the steps).
Rung~3 is different in kind: Palace authors \emph{no} loop, so there is nothing to
render; the obstruction exists only as a marker at the library boundary, and the
hop erases the marker. The distinction is not pedantry---it changes how the
obstruction can ever be discharged. A rendered recurrence (rungs~1--2) records a
loop Palace chose to surface and could, in principle, be re-described; an
opaque-library marker (rung~3) records a boundary Palace never sees inside and that
will not be discharged unless the upstream architecture changes. We will see all
three rungs concretely in Part~IV: the single loop in \cref{ch:krylov}, the nested
double loop in \cref{ch:multigrid}, and the opaque eigensolver in
\cref{ch:eigenvalues}.

\section{The payoff, and the bridge to the machinery}
\label{sec:payoff}

Step back and see what the layered view has bought. By describing the simulator at
a stack of layers rather than once, and by making every move between layers a
rotation with a stated test and a stated justification, we get three things at
once. We get \emph{auditability}: a doubtful reader can check each short move
instead of swallowing one long leap from C++ to calculus. We get \emph{honesty}:
where the abstraction meets a wall, the obstruction is recorded, not papered over.
And we get a \emph{target}: the top layer $L_4$, the calculus you have learned, is
precisely the whole-tensor, no-aliasing form an accelerator backend wants to
consume---the implementation view promised back in \cref{ch:bigidea}.

That is the close of Part~III. You now hold the entire framework: the calculus
(\cref{ch:language,ch:records,ch:tensors,ch:shapes}), the fold and combinators
(\cref{ch:fold,ch:combinators}), the state monad and strata
(\cref{ch:monad,ch:strata}), operators as values (\cref{ch:operators}), and---this
chapter---the map of how all five layers relate. What remains is to watch the
rotations turn on \emph{real algorithms}.

\Cref{part:machinery} does exactly that. Its chapters are where the
explicit-loop view and the combinator view stop being a methodology and become
concrete machine. \Cref{ch:krylov} takes the conjugate-gradient and GMRES solvers
and shows their inner loops rotating from hand-threaded buffers up to
\code{krylov\_step} and \code{iterate\_while}---a single-loop erasure (rung~1 of
\cref{fig:erasure}) with a state-hiding rotation you can now name on sight.
\Cref{ch:multigrid} does the same for the multigrid preconditioner---the very
\code{apply} slot of \cref{wex:slot}---and meets the nested double loop
(rung~2). \Cref{ch:eigenvalues} reaches the eigensolver and meets the opaque
library obstruction (rung~3) head on. Each is a rotation, or an obstruction, of a
kind you have now learned to recognize. Turn the page, and watch them turn.

\summarysection
The simulator is described not once but at a stack of five \emph{layers}: $L_0$,
the cited source; $L_1$, pure functions over immutable values; $L_2$, algebraic
decompositions with performance tricks erased; $L_3$, global tensor-field
operations; and $L_4$, the calculus of Part~III. Each layer is complete and
correct in itself, and adjacent layers are joined by a \emph{rotation}---a
re-expression that turns exactly one impedance (mutation, fusion, iteration, or
coordination) while computing the same numbers. A rotation is a translation across
vocabularies, never a one-to-one renaming. Three tests decide whether a genuine
rotation occurred: it \emph{hides state} (T1), it admits a \emph{different
algorithm in the same slot} (T2, the algorithmic-substitution test), or it
\emph{compresses the threaded state} (T3); if all three fail, the move is a
renaming, a smell to be merged or redesigned away. Every rotation carries one of
five justifications---\code{algebraic}, \code{structural}, \code{reduction\_chain},
\code{empirical\_match} (preferred, because it runs both layers and checks the
numbers), or \code{obstruction}, the first-class negative result when a rotation
honestly does not go through. Obstructions---the sequential reductions and
opaque-library loops that will not lift to whole-field operations---are findings,
not failures, and their \emph{erasure scope} (single loop, nested loop, opaque
call) records how much iteration view a hop erases. The payoff is auditability,
honesty, and an accelerator-ready top layer. Part~IV now turns these rotations on
real numerical algorithms.

\furtherreading
The layered, rotation-based method this book documents is purpose-built, but its
two pillars are standard. The shift from mutable loops to pure, composable,
substitutable pieces is the functional-programming tradition (\cref{ch:bigidea}'s
further reading). The numerical algorithms whose rotations Part~IV makes
concrete---Krylov subspace solvers, multigrid, eigenvalue iterations---are
developed in Saad's \emph{Iterative Methods for Sparse Linear
Systems}~\citep{saad2003}, the standard reference for the solvers behind the
single-loop and nested-loop erasures, and in Trefethen and Bau's \emph{Numerical
Linear Algebra}~\citep{trefethenbau1997}, whose lecture-by-lecture development of
orthogonalization, conditioning, and the eigenvalue problem is the ideal companion
for seeing \emph{why} certain steps resist the iteration rotation. Read either with
the three tests of \cref{sec:tests} in hand, and the place where each algorithm's
loop becomes (or refuses to become) a whole-field operation will jump off the page.

\exercisessection
\begin{enumerate}[nosep]
  \item Name the impedance each of the four rotations in \cref{fig:stack} turns,
  in one phrase each. Which two layers does the \emph{iteration rotation} connect,
  and what is the one situation (\cref{sec:obstruction}) in which it refuses to go
  through?
  \item Classify each move as a genuine rotation or a renaming, and name the
  test(s) it passes or fails: (a) replacing an in-place
  \lstinline[style=l4style]|x.Add(alpha, y)| with a pure
  \lstinline[style=l4style]|axpy alpha y x|; (b) renaming the loop body
  \lstinline[style=l4style]|s += x[i]*y[i]| to a primitive \code{dot} with the same
  sequential accumulation; (c) packaging a preconditioner behind an \code{apply}
  interface that a Jacobi smoother \emph{or} a multigrid V-cycle could implement.
  \item The algorithmic-substitution test (T2) asks whether a \emph{different
  algorithm} could fit a slot. Explain why pasting the name \code{nrm2} onto the
  expression $\sqrt{\sum_i |x_i|^2}$ does \emph{not} pass T2, while wrapping a
  Krylov inner step as \code{cg\_step} (substitutable by \code{minres\_step})
  does. What is the difference between a \emph{name} and an \emph{interface}?
  \item Give the justification kind you would prefer for each rotation and say why:
  (a) unfolding $a x + b y$ into two scaled sums; (b) claiming a hand-rolled GMRES
  matches a reference implementation; (c) a chain of four small associativity and
  linearity steps. When is \code{empirical\_match} unavailable, forcing you back to
  an algebraic argument?
  \item A Gauss--Seidel sweep updates $y_i$ using the \emph{already-updated}
  $y_{i-1}$ (\cref{fig:obstruction}). Explain in one paragraph why this blocks the
  $L_2 \to L_3$ iteration rotation, and why recording the obstruction is more
  useful than either forcing a whole-field form or omitting the operator entirely.
  Which \emph{part} of the operator (body or loop) still rotates cleanly?
  \item Place each on the erasure-scope ladder (\cref{fig:erasure}) and justify the
  rung: (a) a conjugate-gradient solve's outer iteration; (b) a degree-3 Chebyshev
  polynomial smoother with an inner recurrence inside an outer Richardson sweep;
  (c) a call to SLEPc's eigensolver. For which one does Palace author \emph{no}
  loop at all, and what does that imply about whether the obstruction can ever be
  discharged?
  \item \emph{(Synthesis.)} Pick any operator you met in Part~II---a weak-form
  assembly, a boundary-condition elimination, a reduction to a physical
  quantity---and sketch its journey up the stack: what does its mutation rotation
  ($L_0\!\to\!L_1$) hide, does its iteration rotation ($L_2\!\to\!L_3$) succeed or
  hit an obstruction, and what combinator (\cref{ch:fold,ch:combinators}) names its
  coordination at $L_4$? You need not be exhaustive; the point is to walk one
  computation through all five costumes.
\end{enumerate}
